%----------------------------------------------------------------------------------------
% NOTATIE EN FIGUREN VAN DE CURSUS RIJEN
% Staat in de preambule en niet in een hoofdstuk: ook de lijsten met oplossingen en de
% bewaarde grafieken (figuren-cache/) gebruiken dit.
%----------------------------------------------------------------------------------------
% Bewaarde grafiek, geplaatst door maak-pdf.py: bestand en pagina in figuren-cache/
\newcommand{\grafiekbestand}[2]{\includegraphics[page=#2,trim=3bp 3bp 3bp 3bp]{figuren-cache/#1}}

% Kleine grafiek van een rij: \minigrafiek{titel}{ymin}{ymax}{ticks op de y-as}{punten n/u_n}
\newcommand{\minigrafiek}[5]{%
  \pgfmathsetmacro{\eenheid}{2.3/(#3-(#2))}%
  \begin{tikzpicture}[x=.38cm,y=\eenheid cm,>={Stealth[length=1.4mm]},font=\scriptsize]
    \draw[->,lijngrijs] (-.2,0) -- (6.8,0) node[right,text=black] {$n$};
    \draw[->,lijngrijs] (0,#2) -- (0,#3) node[above,text=black] {$u_n$};
    \foreach \x in {1,...,6} \draw[lijngrijs] (\x,0) ++(0,1.2pt) -- ++(0,-2.4pt) node[below=-1pt,text=grijs] {\x};
    \foreach \y in {#4} \draw[lijngrijs] (0,\y) ++(1.2pt,0) -- ++(-2.4pt,0) node[left=-1pt,text=grijs] {$\y$};
    \foreach \x/\y in {#5} \fill[accent] (\x,\y) circle (1.6pt);
    \node[anchor=south,font=\sffamily\bfseries\footnotesize] at (3.4,#3) {#1};
  \end{tikzpicture}}

% Tekst of berekening links, figuur rechts: \naast{links}{rechts}
\newsavebox{\naastbox}
\newcommand{\naast}[3][]{%
  \par\noindent\sbox{\naastbox}{#3}%
  \begin{minipage}[t]{\dimexpr\linewidth-\wd\naastbox-3mm\relax}\vspace{0pt}#2\end{minipage}\hfill
  \begin{minipage}[t]{\wd\naastbox}\vspace{0pt}\usebox{\naastbox}\end{minipage}\par\addvspace{1mm}}
% Een rij figuren, gecentreerd
\newcommand{\figuren}[1]{\par\addvspace{1mm}{\centering #1\par}\addvspace{1.5mm}}
% Opgave met een grafiek rechts: \begin{metgrafiek}{tikz} opgave en deelvragen \end{metgrafiek}
\newsavebox{\grafiekbox}
\newenvironment{metgrafiek}[2][]{%
  \sbox{\grafiekbox}{#2}%
  \begin{minipage}[t]{\dimexpr\linewidth-\wd\grafiekbox-3mm\relax}%
}{%
  \end{minipage}\hfill
  \begin{minipage}[t]{\wd\grafiekbox}\vspace{-2mm}\usebox{\grafiekbox}\end{minipage}\par}

% Zoals stappen, met eigen kolombreedtes: \begin{stappenvrij}{6.2cm}{3.8cm}. Voor naast een figuur.
\newenvironment{stappenvrij}[2]{%
  \par\renewcommand{\arraystretch}{1.25}%
  \noindent\begin{tabular}{@{}>{$\displaystyle}p{#1}<{$}@{\hspace{3mm}}>{\RaggedRight\small\itshape\arraybackslash}p{#2}@{}}%
}{\end{tabular}\par}
